\chapter{Introducción}
\label{ch:intro}

Litoteca es un proyecto de visión por computador aplicado al reconocimiento de
muestras de roca. Su contexto es educativo y de divulgación: una persona puede
encontrar una piedra en campo, en una colección didáctica o en un ejercicio de
clase, y necesitar una primera orientación sobre qué tipo de roca está observando.
El área de aplicación es el reconocimiento de objetos mediante clasificación de
imágenes, donde una fotografía tomada con el celular o cargada desde el
dispositivo se analiza para predecir la clase más probable entre nueve tipos de
roca.

El problema no se limita a asignar una etiqueta. Clases como chert, diatomite o
siliceous-sinter pueden compartir rasgos visuales, por lo que el sistema debe
acompañar la predicción con información interpretable: confianza, alternativas
probables y una ficha que explique qué es la roca, cómo se formó y con qué puede
confundirse. La predicción funciona como apoyo visual y pedagógico, no como
sustitución de un análisis de laboratorio.

El objetivo del trabajo es implementar un modelo de clasificación de rocas en
Azure Custom Vision e integrarlo en una aplicación web práctica. El prototipo
cubre nueve clases (basalto, basalto olivínico, sílex/chert, arcilla,
conglomerado, lutita, diatomita, sinter silíceo y yeso), usa un dataset preparado
y depurado, y se despliega como sitio estático en Cloudflare Pages con una Pages
Function que actúa como proxy de predicción. Al cierre de esta versión, la
aplicación y el flujo de despliegue están listos y ya se cuenta con una primera
iteración Quick entrenada (precisión 83.08\%, recall 52.43\%, AP 72.74\%), usable
para demo mientras se evalúan mejoras con iteraciones adicionales.
